%% kapa: exam=1 round=31 year=2020 period=1 form=A subject=civil_law subject_ko=민법 seq=32 number=32
%% source: Q-Net 감정평가사 기출문제 (cst00309) artlSeq=5207219 file=extracted/1008/5207219/A형/A형/2020년도 제31회 감정평가사 1차 시험_1교시_A형.pdf
%% conversion: kordoc 4.24.51 -> tools/gen1.py (auto); review=proofread-sonnet; flags=none
\begin{question}{32}
명의신탁에 관한 설명으로 옳지 않은 것은? (다툼이 있으면 판례에 따름)
\begin{choices}
  \item 종중재산이 여러 사람에게 명의신탁된 경우, 그 수탁자들 상호간에는 형식상 공유관계가 성립한다.
  \item 3자간 등기명의신탁관계의 명의신탁자는 수탁자에게 명의신탁된 부동산의 소유권이전등기를 청구하지 못한다.
  \item 채무자는 채권을 담보하기 위하여 채권자에게 그 소유의 부동산에 관한 소유권이전등기를 할 수 있다.
  \item 부부 사이에 유효하게 성립한 명의신탁은 배우자 일방의 사망으로 잔존배우자와 사망한 배우자의 상속인에게 효력을 잃는다.
  \item 계약 상대방이 명의수탁자임을 알면서 체결한 매매계약의 효력으로 소유권이전등기를 받은 사람은 소유권을 취득한다.
\end{choices}
\end{question}
\answer{4}
